\subsection{New reductions for the unrestricted sixteen-arrow baseline}

\begin{lemma}[Proved: bounded act cones through monoid order eight]
\label{CP2:A-baseline-eight-act}
Let $M$ be a finite monoid of order at most eight and $S$ a
nonempty finite right $M$-set.  If the act-cone category has
bounded rational homology, all its rational degree-two products
are zero.  The left-act assertion is dual.  The same holds
after adjoining an opposite endpoint with projective trivial
rational module on the appropriate side.
\end{lemma}
\begin{proof}
Put $R=\mathbb QM$ and
$L=\ker(\mathbb QS\xrightarrow{\epsilon}\mathbb Q)$,
with its actual right $R$-module structure.  To specify the
boundedness implication precisely, take a projective left
resolution $P_*\to\mathbb Q$.  The action-space map to
$BM$ is computed by
\[
 \mathbb QS\otimes_R P_*\longrightarrow
 \mathbb Q\otimes_R P_*.
\]
Projective left modules are flat, so
\[
 0\to L\otimes_RP_*\to\mathbb QS\otimes_RP_*
   \to\mathbb Q\otimes_RP_*\to0
\]
is exact in every degree.  The mapping cone of its last
map is quasi-isomorphic to $(L\otimes_RP_*)[1]$.
The two-block description of the act cone therefore gives
\[
 \widetilde H_j(B\mathcal C(M,S);\mathbb Q)
   =\operatorname{Tor}_{j-1}^R(L,\mathbb Q)
   \qquad(j\geq1).
\]
Thus boundedness means eventual vanishing of these particular
Tor groups.  It is not inferred separately for $BM$ or its
action space, whose homology can indeed cancel in a cofiber.

If $H^2(BM;\mathbb Q)=0$, second classes of the cone lift
from its suspension quotient and all their products vanish.
If there are at most two minimal rows, use
Lemma~\ref{CP2:A-high-two-row-act}.  Otherwise nonzero
vertex second cohomology requires at least two columns,
while there are at least three rows.  The minimal-ideal
size and the identity outside it force a trivial maximal
group, exactly a $3\times2$ band $K$, and at most two
outside elements.  Put $A=R/\mathbb QK$.

The ideal $\mathbb QK=Re\otimes_{\mathbb Q}eR$ is
projective on the left, so $\operatorname{pd}_R A\leq1$.
The row augmentation is an exact sequence
\[
 0\to V\to Re\to\mathbb Q\to0,
\]
where $KV=0$ and $\dim V=2$.  Hence $V$ is a left
$A$-module and, for $j\geq2$,
$\operatorname{Tor}_{j}^R(L,\mathbb Q)
=\operatorname{Tor}_{j-1}^R(L,V)$.

A unital rational algebra of dimension at most two is
semisimple or the dual-number algebra
$\mathbb Q[\varepsilon]/(\varepsilon^2)$.
In the semisimple case $V$ is a summand of a finite
sum of copies of $A$, so it has $R$-projective dimension
at most one and $\operatorname{Tor}_3^R(L,\mathbb Q)=0$.
In the dual-number case Jordan decomposition gives
$V=A^a\oplus k^b$, where $k=A/(\varepsilon)$.
Restrict the exact sequence
$0\to k\xrightarrow{1\mapsto\varepsilon}A\to k\to0$
to $R$.  Since $\operatorname{Tor}_{j}^R(L,A)=0$ for
$j\geq2$, its long exact sequence gives
\[
 \operatorname{Tor}_{j}^R(L,k)
   \cong\operatorname{Tor}_{j-1}^R(L,k)
   \qquad(j\geq3).
\]
If $b=0$, the projective-dimension argument already
applies.  If $b>0$, eventual vanishing of the cone
homology implies eventual vanishing for $V$, hence
for $k$.  The displayed recurrence then forces
$\operatorname{Tor}_2^R(L,k)=0$.  In either case
\[
 H_4(B\mathcal C(M,S);\mathbb Q)
   =\operatorname{Tor}_3^R(L,\mathbb Q)
   =\operatorname{Tor}_2^R(L,V)=0.
\]
This proves the obstruction in the remaining case.
Opposition gives left acts, and the actual exact-orbit
functor gives the projective opposite-endpoint version.
\end{proof}

\begin{proposition}[Proved: large corners in a linear sixteen-arrow profile]
\label{CP2:A-baseline-linear-large}
A three-object category with full linear acyclic support and
at most sixteen arrows cannot realize $\mathbb{CP}^2$ if any
corner has order at least eight.
\end{proposition}
\begin{proof}
Write the corner orders as $a,b,c$, and the hom-set sizes as
$x=|C(0,1)|$, $y=|C(1,2)|$, $z=|C(0,2)|$; all six
numbers are positive.  First suppose $c\geq8$.  Then
$a+b+x+y+z\leq8$ and $a\leq4$.  If $x,z\leq2$,
the small-source binary-stalk lemma excludes the square.
Otherwise $x\geq3$ or $z\geq3$.

If $x\geq3$, at least one of $y,z$ equals one: having
both at least two would already give
$a+b+x+y+z\geq9$.  If both are one, the sink cut is
a full singleton join.  If exactly one is one, the lower
two-object block has at most five arrows, so its rational
second cohomology is zero by the H2-cost lemma.  The
singleton incoming arrow to the sink then excludes all
second products by the fixed-arrow corollary.

If $z\geq3$ and $y=1$, the lower block has at most four
arrows and the same fixed-arrow argument applies.  If
$y\geq2$, the inequality forces
$a=b=x=1$, $y=2$, $z=3$, and $c=8$.  The comma category
for deleting the middle object has its one object given
by the unique arrow $0\to1$ and has trivial endomorphism
monoid.  Quillen's Theorem~A identifies the whole category
with the ordinary three-state act cone over the order-eight
sink.  Its boundedness follows from this equivalence, so
Lemma~\ref{CP2:A-baseline-eight-act} excludes its square.
Opposition handles a large source corner.

Finally suppose $b\geq8$.  Both endpoint corners have
order at most four.  Either $x,z\leq2$ or $y,z\leq2$
gives the binary-stalk obstruction.  Avoiding both would
require $z\geq3$ or both $x,y\geq3$.  The latter costs
at least seventeen arrows, so $z\geq3$ and
$a+c+x+y\leq5$.  In particular $a+c\leq3$, and at
least one of $x,y$ is one.  If $x=1$ but $a>1$,
then the same inequality forces $c=y=1$; the dual
statement holds if $y=1$ but $c>1$.  Thus there is
always a trivial endpoint with its adjacent hom-set
singleton.  Its singleton comma category permits
deleting the middle object.  The remaining two corners
have orders at most two, hence semisimple rational
algebras; the acyclic edge has rational homology only
through degree one.  Its fourth homology is zero.
All large-corner positions are excluded.
\end{proof}

\begin{proposition}[Proved: four-object large corners through sixteen arrows]
\label{CP2:A-baseline-four-large}
A four-object category with at most sixteen arrows and a
corner of order at least eight cannot realize $\mathbb{CP}^2$.
Consequently, to establish the four-object portion of the
unrestricted sixteen-arrow baseline it suffices to consider
acyclic support, exactly sixteen arrows, and corners of
order at most seven.
\end{proposition}
\begin{proof}
If the classifying space were $\mathbb{CP}^2$, the
excess-eight theorem would force $n=16$, and
Corollary~\ref{CP2:A-high-nine-support} would give
acyclic support and no object retract.  Let $v$ have
corner $M$ of order $q\geq8$.  With $d$ the corner
excess away from $v$, the budget becomes
\[
 d+P+R=10-q\leq2.
\]
The case $P=0$ is excluded by singleton-support
comparison, because a four-object poset has no
four-dimensional homology.

If $v$ is nonextreme, then $R\geq1$.  The remaining
budget forces $q=8$, $d=0$, $P=R=1$.  There is one
triangle through $v$ and one leaf edge; $v$ is its
middle object and has no other neighbor, since any
such comparison would create another cycle.
One adjacent hom-set is singleton.  Its other
endpoint is trivial and gives a singleton comma
category for deleting $v$, on the appropriate
side of Quillen's Theorem~A.  The result has three
trivial corners and acyclic support, hence nerve
dimension at most two.  This contradicts the
nonzero fourth homology.

Now take opposites so that $v$ is a sink, and let
$D=C\setminus\{v\}$.  If it has a singleton incoming
hom-set, then $|D|\leq7$.  At equality $q=8$ and
there is exactly one incoming arrow, so $D$ is
connected.  The H2-cost lemma gives $H^2(BD;\mathbb Q)=0$
in that case.  When $|D|\leq6$, each connected
component of size two has at most five arrows and
each isolated corner has order at most four;
the same lemma and the small-monoid result give
$H^2(BD;\mathbb Q)=0$, even if $D$ is disconnected.
The singleton is a fixed incoming arrow, and
Corollary~\ref{CP2:A-fixed-arrow} applies.

Suppose every incoming hom-set has at least two
elements.  There cannot be three neighbors, which
would require at least $8+3+6=17$ arrows.
With two neighbors the budget forces $q=8$, all
other corners trivial, two doubled incoming
hom-sets, and one additional singleton edge
connecting the fourth object $w$ to a neighbor
$x$.  This edge is $x\to w$: the reverse direction
would create a third hom-set into $v$.  Delete
the trivial leaf $w$.  The remaining star has
two trivial leaves and the order-eight center.
Identifying its two leaf points gives an ordinary
left-act cone with four states.  On spaces this
adds only a circle, so preserves bounded rational
homology and cohomology products in degrees at
least two.  Lemma~\ref{CP2:A-baseline-eight-act}
therefore kills its square.

With one neighbor $w$, write $s=|C(w,v)|\geq2$.
There is no incoming cross arrow to $w$ in $D$,
since its composite would create another neighbor
of $v$.  The nerves of $D$ and the edge collage
$Y$ on $w,v$ cover $NC$, meeting in the corner
at $w$.  The category $D$ is connected and needs
at least five arrows, so $s\leq3$.  If $s=2$,
$|D|\leq6$ and its second cohomology vanishes;
the shared corner has order at most two and is
rationally acyclic.  Fourth restriction is
injective, and the universal two-cross theorem
kills the square on $Y$.  If $s=3$, equality
forces $q=8$ and $|D|=5$: $D$ is a trivial-corner
singleton-hom tree and is contractible.  The
pushout identifies $BC$ with the three-state
ordinary act cone $BY$.  The order-eight act
lemma again excludes its square.  These cases
exhaust the budget.
\end{proof}

\begin{lemma}[Proved: a five-element extreme beyond two trivial corners]
\label{CP2:A-baseline-trivial-pair-five}
Suppose a three-object category has full linear acyclic support,
its first two corners are trivial, and its last corner has
order at most five.  Its rational degree-two products are
zero, with no bound on its hom-set sizes.  The dual statement
holds at the other end.
\end{lemma}
\begin{proof}
The full subcategory $D$ on the first two objects has graph
classifying space.  If the last corner $M$ is not $P_5$,
its second rational cohomology is zero by the small-monoid
classification, and the two-block lemma applies.

For $M=P_5$, the category of elements of the incoming
$D$-act is itself a graph $\Gamma$: it has one vertex for
each incoming arrow from either source object, and one edge
for each pair consisting of an arrow of $D$ and an incoming
arrow from its target.  There are no composable pairs of
nonidentity arrows inside $D$, so there are no higher
nondegenerate simplices.  The commuting $P_5$ action sends
vertices to vertices and each parametrized edge to the edge
of its image pair.  Thus it is precisely a finite graph
action of the kind in Lemma~\ref{CP2:A-high-P5-graph},
after taking opposites if needed; $P_5$ is isomorphic to
its opposite by transposing its band coordinates.

The two-block model identifies the quotient which collapses
$BD$ with the cofiber of $\Gamma_{hP_5}\to BP_5$.
Since $H^2(BD;\mathbb Q)=0$, the quotient map induces a
surjection onto $H^2(BC;\mathbb Q)$ by the relative
cohomology sequence.  Every second class is therefore the
pullback of a class whose products vanish by the graph
cofiber lemma.  Naturality of the cup product proves the
claim.  Neither boundedness nor a rank-one assumption was
used.
\end{proof}
